\section{The limits of regional identification}
\label{sec:robustness}

The obvious way to move from the accounting of Section~\ref{sec:results} to a
causal statement about prices is to exploit the autonomous regions, which may
set a minimum wage above the mainland's. This section reports why that design
cannot be carried in Portugal. We present it as a result, because the reasons
are specific and checkable rather than a general complaint about power.

\subsection{How much variation actually exists}

Under a specification with region-by-category and calendar-time fixed effects,
what identifies a coefficient is the set of region-months in which the statutory
change differs from the mainland's. Counting them is not the same as counting
observations, and the difference is large.

The Azores contribute nothing. Their supplement is proportional, so
$\Delta \log MW$ in the Azores equals $\Delta \log MW$ on the mainland exactly
in every month after the supplement took effect; the region adds rows to the
panel and not one degree of freedom.

Madeira is more interesting, and the reason is a fact about its law rather than
about our data. Madeira legislates a euro figure rather than a rule, which looks
like independent variation, but until 2016 the figure was not independent: the
preamble to Decreto Legislativo Regional n.º~18/2016/M states the regional
policy as a \MadeiraRuleFactorPct{}~per cent supplement on the mainland amount,
begun in 1987, and the acts obey it to the cent from \MadeiraRuleFirstYear{} to
\MadeiraRuleLastYear{}: over those years Madeira's figure is the mainland's
multiplied by that supplement, with nothing left over. Madeira was on the same
kind of rule as the Azores, written out by hand year by year, and contributed no
more timing than they did. After \MadeiraRuleLastYear{} the acts leave the rule
and the premium widens, reaching \MadeiraPremiumEndPct{}~per cent by
\MadeiraPremiumEndYear{}. Madeira's usable
variation therefore begins in 2017, not in 1987, and no amount of further
retrieval changes that.

Two features of the law shape what is counted. A regional decree stays in force
until another replaces it --- DLR~18/2016/M expressly revokes DLR~13/2014/M ---
so a year without an act is a year the previous figure continued, not a year
without a wage. And because Madeira states a figure rather than a rule, a
standing figure can be overtaken: its 2017 act set 570 euros, which was still
the region's own last word in 2018 when the national wage reached 580. The
national wage is a floor, so 580 is what bound, and Madeira's premium was
extinguished for a year without any regional act saying so. Both are read from
the acts and neither is inferred from a summary table.

Over the estimation window \IdentifyingEvents{} region-months diverge from the
mainland, across \IdentifyingRegions{} region.

A gap in the register would still not be neutral: if a region legislated in a
year whose act is missing, its wage would appear frozen against a rising
national one, which is indistinguishable from a deliberate freeze and enters the
estimation as a policy shock that never happened. The pipeline therefore
requires known-incomplete years to be declared and defers to the national wage
in them. Madeira's acts are registered from 2010; the years before that are
declared, and the window starts where the register is contiguous.

\subsection{Why conventional inference cannot be trusted here}

Cluster-robust standard errors are justified asymptotically in the number of
clusters, and Portugal has few regions. That conventional inference over-rejects
in this situation, and that a wild cluster bootstrap corrects it, is established
\citep{cameron2008bootstrap}. In simulation with seven clusters and a
true null, the conventional test rejects at the five per cent level about twenty
per cent of the time, while a restricted wild cluster bootstrap holds its size.
The estimation panel below has \BootstrapClusters{} clusters, and at that count
the procedure \BootstrapBasis{}.

This is not a refinement. It is the difference between a result and an artefact,
and Table~\ref{tab:regional-design} shows it on the actual data.

\begin{table}[htbp]
  \centering
  \caption{Cumulative price response to the statutory change, regional design.
  Estimated over the window in which Madeira's register is contiguous, with
  region-category and calendar-time fixed effects, clustered on region. The
  interval is obtained by inverting the bootstrap test rather than from the
  standard error, so it and the $p$-value beside it are the same procedure read
  two ways; an asterisk would mark an interval whose search did not close, and
  \IntervalsBounded{} of these closed.}
  \label{tab:regional-design}
  \input{tables/regional_design}
\end{table}

Conventional inference calls \ClusteredHighlySignificant{} of the
\HorizonsEstimated{} horizons significant at the one-tenth of one per cent
level. The bootstrap leaves \BootstrapSurvivors{} of them standing, at
$p =\ $\MinBootstrapP{}. Had we reported only the third column, the paper would
have claimed substantial pass-through of minimum-wage costs into consumer prices
at nearly every horizon.

The last column is the one to read, and it is the column a $p$-value cannot
supply. At \PeakHorizon{}~month the interval runs from \PeakIntervalLower{} to
\PeakIntervalUpper{}, so what the data exclude is a response above
\PeakIntervalUpper{} and below \PeakIntervalLower{}, which leaves everything
from a negligible pass-through to one several times the point estimate. At
\WidestIntervalHorizon{}~months it runs from \WidestIntervalLower{} to
\WidestIntervalUpper{} and excludes almost nothing of interest. These are not
standard-error bands: each is the set of values the reported bootstrap test does
not reject, obtained by inverting it, so an interval excluding zero and a
$p$-value below five per cent are guaranteed to agree. A band built the
conventional way, by resampling around the estimate, did not agree --- it
declared horizons significant that the null-imposed test does not reject --- and
that disagreement is why it was discarded.

That surviving horizon is not a result, and the reason is the same one that
motivated the bootstrap. \HorizonsEstimated{} horizons of one series are a
family of tests, not \HorizonsEstimated{} papers. Reporting the smallest of
several $p$-values as though it stood alone rejects a true null far above its
nominal size --- with seven tests at five per cent, the chance of at least one
rejection approaches a third --- which is the error the bootstrap was introduced
to correct, one level up. Holm's step-down adjustment across the horizon family
controls it under arbitrary dependence between horizons, which is what the
strong serial correlation of a price series requires. Adjusted, the smallest
$p$-value is \MinBootstrapHolmP{} and \BootstrapSurvivorsHolm{} horizons
survive.

What has changed is the plausibility of the point estimates, not the strength of
the evidence. The largest coefficient is now \PeakCoefficient{} at
\PeakHorizon{}~month, where an earlier version of this design --- estimated on a
window in which the register held three identifying events --- returned a peak
above one, implying prices rising by more than the wage shock that caused them.
The profile is now the shape a pass-through function ought to have: positive and
modest on impact, decaying within the year. We still decline to interpret it. It
is identified by \IdentifyingEvents{} events in a single small, remote,
tourism-intensive region, so anything else moving Madeiran prices on those dates
enters the coefficient with nothing to separate it out.

\subsection{Do the leads predict prices before the policy moves?}

The falsification battery this project set itself requires that leads of the
statutory change carry no information about pre-treatment inflation. We test the
\PreTrendLeads{} leads jointly rather than reading them off a plot, because
inspecting them one at a time both multiplies the chance that one looks
significant and misses a trend spread thinly across several without any single
one standing out.

The joint restriction \PreTrendRejects{} the null at five per cent, with a
bootstrap $p$-value of \PreTrendP{}. The gap between references is the more
instructive number. The Wald statistic is \PreTrendWald{}, which against the
asymptotic $\chi^2$ reference would be overwhelming evidence of a pre-trend;
referred instead to its own wild cluster bootstrap distribution, it sits at a
tenth. With nine clusters the asymptotic reference is not usable, and this is
the clearest illustration of that fact anywhere in the paper.

Two cautions attach to reading this as a pass. A design that struggles to detect
a treatment effect will also struggle to detect a pre-trend, so failing to
reject is weak evidence of the absence of one. And $p=\PreTrendP{}$ is not
comfortable: it is the kind of value that would become a rejection under a
modestly different window or lead count. We report it as a diagnostic that does
not disqualify the design, not as a clean bill of health.

\subsection{The exposure design, built, and why it is too flat}

The standard substitute for regional policy variation is regional exposure to a
common national change: a region whose employment sits disproportionately in
low-paying industries faces a larger effective shock. Writing $q_{rs,0}$ for
region $r$'s baseline employment share in industry $s$ and $b_{s,0}$ for the
national share of employees paid the minimum wage in that industry, exposure is
$B_{r} = \sum_{s} q_{rs,0}\, b_{s,0}$ \citep{goldsmithpinkham2020bartik}.

Both terms are available, and we report the measure rather than assert that it
cannot be had. Eurostat's regional accounts give $q_{rs,0}$ for the
\ExposureRegions{} NUTS~II regions and the monitoring reports of
Section~\ref{sec:data} give $b_{s,0}$ nationally. The regional accounts publish
coarser activity groups than the bite is measured at, so each group's bite is
the mean of its sections weighted by national employment in them; taken
unweighted the mean puts manufacturing on equal footing with a near-empty
utilities sector, and reverses the regional ordering.

Frozen on 2015 composition and the October~2015 survey round, exposure takes \ExposureDistinctValues{} distinct
values across the \ExposureRegions{} regions, running from \ExposureMinPct{} to
\ExposureMaxPct{}~per cent. That is the whole difficulty. The high-low spread is
\ExposureSpreadPP{}~percentage points and the coefficient of variation is
\ExposureCV{}. Regional industry composition varies
a great deal in Portugal --- the Algarve holds nearly twice the Alentejo's share
of employees in trade, transport and hospitality, and the Alentejo holds more
than twenty-five times Grande Lisboa's share in agriculture --- but the bite
does not differ enough between industries for composition to translate that into much dispersion
in exposure. A measure that differs across regions establishes that a regional
effect is identified in principle; it does not establish that the effect is
identified precisely enough to be informative, and this spread does not.

We can now do better than argue from the spread, because the published table
reports five survey rounds rather than one. Its October~2015 round is
predetermined for every statutory change from 2016 onwards, so the design can be
estimated rather than merely described. Table~\ref{tab:exposure-design} reports
it.

\begin{table}[htbp]
  \centering
  \caption{Cumulative price response to the statutory change interacted with
  predetermined regional exposure. Exposure is frozen on 2015 composition and
  the October~2015 survey round; the window begins in 2016, so exposure precedes
  every shock. Region-category and calendar-time fixed effects, clustered on
  region. The interval inverts the bootstrap test, as in
  Table~\ref{tab:regional-design}; \ExposureIntervalsBounded{} of these closed
  inside the search.}
  \label{tab:exposure-design}
  \input{tables/exposure_design}
\end{table}

This design has an advantage the previous two lack. Because exposure varies
across regions, the interaction survives calendar-time fixed effects, so
everything moving national prices in a month --- energy, taxes, the pandemic ---
is absorbed rather than left in the error. It is the specification the
literature favours, and here it is correctly identified.

It rejects nothing. Not one of the \ExposureHorizons{} horizons is significant
at five per cent even by conventional clustered inference, whose smallest
$p$-value is \ExposureMinClusteredP{}; the bootstrap and the Holm correction
have nothing left to reject. Inference known to over-reject cannot manufacture a
single false positive here, and the coefficients wander as far as
\ExposureWidestCoefficient{} at \ExposureWidestHorizon{}~months without
approaching significance.

That is the weakest possible way to put it, and the intervals say why. At
\ExposureWidestIntervalHorizon{}~months the inverted interval runs from
\ExposureWidestIntervalLower{} to \ExposureWidestIntervalUpper{}. A coefficient
of one would mean the wage rise passed into prices one for one in the most
exposed region relative to the least; this interval admits that, admits sixty
times that, and admits the same magnitudes with the sign reversed. Reporting
only that the design fails to reject would invite the reading that it found no
pass-through. It did not find no pass-through. It failed to locate the
coefficient anywhere within two orders of magnitude, and no null result stated
as a $p$-value can convey that.

The same falsification test applies here, and this design passes it more
comfortably than the policy design does. The \ExposurePreTrendLeads{} leads
jointly \ExposurePreTrendRejects{} the null, at $p=\ $\ExposurePreTrendP{}.
That is reassuring only up to a point: a design this imprecise would struggle to
detect a pre-trend as surely as it struggles to detect an effect.

The null does not depend on the choices we made in building the measure.
Table~\ref{tab:exposure-robustness} re-estimates it varying the survey round the
bite is taken from, the year composition is frozen at, the bite assumed for the
sectors the survey misses, and which region is dropped. No horizon in any of the
sixteen specifications survives the Holm correction, and none rejects at five
per cent before it.

\begin{table}[htbp]
  \centering
  \caption{The exposure design under alternative construction choices.
  Coefficient range is across horizons within a specification; rejections are at
  five per cent, before and after the Holm correction across the horizon family.}
  \label{tab:exposure-robustness}
  \input{tables/exposure_robustness}
\end{table}

The coefficients themselves tell the more useful story. They range from
$-2.36$ to $0.49$ in the baseline and from $-9.89$ to $0.47$ when Grande Lisboa
is dropped. Grande Lisboa is the least exposed region by a clear margin, so
removing it collapses most of the contrast the design runs on, and the estimate
moves by an order of magnitude in response. A coefficient that behaves that way
is not measuring a price response; it is dividing by a number close to zero.
That the null survives every specification is therefore not evidence that
pass-through is absent. It is evidence that this design would fail to detect it
either way.

Two further qualifications keep the exposure measure from being read as more
than it is, and both weaken it.

The exposure measure is conditional on the sectors the survey covers. Coverage
runs from \ExposureCoverageMinPct{} to \ExposureCoverageMaxPct{}~per cent of a
region's employees, and the uncovered part is not spread evenly: agriculture is
unsurveyed and is a far larger share of employment in the Alentejo than around
Lisbon. Dividing it out therefore compresses exactly the regional differences
the measure exists to capture. Completing the measure under an assumed bite $u$
for the unsurveyed sectors, rather than conditioning them away, widens the
spread substantially --- at $u=0$ it is roughly seventy per cent wider than the
conditional figure --- though it stays within a few percentage points across
every assumption we regard as defensible. The regional \emph{ranking}, however,
is not stable across those assumptions, so we do not report one.

Nor does failing to reject establish that the effect is zero. What a null is
worth depends on what the design could have detected. For a ten per cent
statutory rise, the difference in price response between the most and least
exposed region is resolved to about four tenths of a percentage point at the
horizons where the estimate is tightest, which would be informative; but at six
and eighteen months the interval spans more than three percentage points, which
is wider than any pass-through estimate in this literature. At those horizons
the finding is not that there is no response. It is that this design cannot see
one. Those figures are the conventional standard-error bands, and they are the
optimistic version: the inverted intervals in
Table~\ref{tab:exposure-design}, which are what the reported test actually
supports, are wider still by a large factor.

What can be said is narrower than "exposure is flat in Portugal". It is that
minimum-wage incidence within the sectors where a comparable bite is observed
varies little once weighted by regional industrial composition, and that the
resulting design is under-powered at most horizons. Publishing the bite by
region would address both; publishing more of the same national bite would
address neither.

\subsection{The category-differential design, and why it also falls short}

The remaining variation is across consumption categories facing a common
national change, with the differential response estimated directly, one
coefficient per category. It is reported alongside the exposure measure because
it does not inherit the assumption that within-industry coverage is constant
across regions.

That design cannot carry calendar-time fixed effects, because a national shock
does not vary across regions within a month; time effects would absorb it
entirely. It therefore identifies each category relative to the others and is
exposed to anything else moving national prices at the same time. The estimated
ranking is partly plausible: food and, next, restaurants and accommodation sit
at the top, and accommodation and food services is the activity with by far the
highest measured minimum-wage coverage. But \SeasonalWorstCategory{} returns a
coefficient an order of magnitude more negative than any other category's, and
diagnosing it defeats the design.

An earlier version of this paper put that coefficient down to a break or a
rebasing in the series. That was wrong, and the true explanation is a property
of the policy rather than of the data. Portugal moves its minimum wage on
1~January: \ShockModalSharePct{}~per cent of all statutory log change over the
sample falls in \ShockModalMonth{}, so the shock is very nearly an indicator for
that one month. \SeasonalWorstCategory{} falls by \SeasonalWorstSwingPct{}~per
cent in an average \ShockModalMonth{}, because winter sales enter the index and
have done so more sharply since the harmonised treatment of seasonal products
took effect in 2011. No other division moves by as much as a point. A
specification with no month-of-year effects therefore reads the January sales
cycle as the price response to the January wage rise, and it is the size of the
sales cycle, not any defect in the series, that makes the coefficient extreme.
Category-specific linear trends do not help, because the confound is seasonal
and not a trend.

Absorbing month-of-year effects removes the artefact and takes the design with
it. Only \ShockSurvivingSharePct{}~per cent of the shock's variation survives
those effects, because most of the shock \emph{is} the calendar. What remains is
the handful of years whose January increase differed in size from the others,
plus the single mid-year change of October 2014. We therefore do not report the
category ranking, and the reason is not that one series is corrupt but that a
policy which always moves in the same month cannot be separated from that month
without discarding the variation that identifies it. The point generalises to
any country that updates its wage floor on a fixed calendar date.

\subsection{The region-by-category design, and the magnitudes that condemn it}

The two designs above fail in different ways, and each failure is what the other
would need to survive. The exposure design cannot rule out that its estimates
are driven by tourism cycles and island supply shocks, because with a shock that
varies only across regions it cannot carry a region-time effect. The
category-differential design is defeated by the January sales cycle, because
with a shock that varies only across categories it cannot carry a category-time
effect. A shock varying on both dimensions can carry both, and that is the case
for building one.

The fuller exposure is structural rather than reduced-form. Writing $\ell_s$ for
labour's share of costs in industry $s$ and $\omega_{cs}$ for the share of
consumption category $c$ supplied by it,

\begin{equation}
  B_{rc} = \sum_{s} q_{rs}\, b_{s}\, \ell_{s}\, \omega_{cs},
  \label{eq:structural-exposure}
\end{equation}

and the regressor is $B_{rc}\,\Delta \log MW^{PT}_t$. An earlier version of this
paper argued that the two new terms could not help, since both are national. That
was a reasoning error rather than a data problem: with $\alpha_{rc}$,
$\lambda_{rt}$ and $\mu_{ct}$ absorbed, what identifies the coefficient is the
non-additive part of $B$, and a national factor enters it through its product
with regional composition. Composition and the bridge are not proportional to one
another, so the product is not additively separable and does not vanish.

Two things had to be built. The labour-cost share comes from national accounts.
The bridge does not exist in any published source --- nothing crosses a
consumption purpose with a producing industry for Portugal --- so it is composed
from household final consumption by product, taken at basic prices and domestic
uses, and a concordance recorded with its reasoning in
\texttt{config/consumption\_bridge.yaml}. Both choices about the source matter
more than the concordance does. At purchasers' prices the retail margin on a
good is credited to the industry that made it rather than the one that sold it,
which moves the wage exposure of the shopping basket between industries; and a
quarter of Portuguese household spending is imported content or product taxes,
with no Portuguese producing industry behind it at all.

What the design has is smaller than the measure suggests. Because the statutory
change is common, the region-time effects remove the row means of $B$ and the
category-time effects its column means, so the identifying variation is the
double-demeaned remainder: \StructuralIdentifyingSpread{}~percentage points
across \StructuralCells{} region-category-months in
\StructuralClusters{} regions. Table~\ref{tab:structural-design} reports it.

\begin{table}[htbp]
  \centering
  \caption{Cumulative price response to the statutory change interacted with
  region-by-category structural exposure, equation~\eqref{eq:structural-exposure}.
  Region-category, region-time and category-time fixed effects, clustered on
  region. The coefficient is on an exposure index rather than a cost share and
  is not interpretable on its own; the text scales it into points.}
  \label{tab:structural-design}
  \input{tables/structural_design}
\end{table}

It is the only design in this paper that rejects anything.
\StructuralBootstrapSignificant{} of the \StructuralHorizons{} horizons are
significant at five per cent by the bootstrap, at
$p=\ $\StructuralMinBootstrapP{}, and \StructuralHolmSignificant{} survive the
Holm correction across the horizon family --- the same correction, for the same
reason, that leaves nothing standing in the regional design.

We would not report those horizons as a finding even had they survived it, and
the reason is a magnitude rather than a $p$-value. The coefficient cannot be
read directly. The exposure weights each industry by the region's employment
share in it, which is what supplies the regional variation and also means it is
not a cost share but a fraction of one; the fraction runs from under a hundredth
to two fifths across cells, so no single rescaling makes the coefficient
interpretable. What is interpretable is the response it implies. Scaled into
points, the estimate at \StructuralPeakHorizon{}~months gives a differential
price response of \StructuralPeakResponsePP{}~percentage points between the most
and least exposed cells for a \StructuralRisePct{}~per cent statutory rise.

That figure can be checked against a ceiling the design itself supplies. The
most minimum-wage-intensive consumption category has
\StructuralCostShareCeilingPct{}~per cent of its costs in minimum-wage labour,
so if every euro of that cost reached prices, a \StructuralRisePct{}~per cent
rise could move it by at most \StructuralPassThroughCeilingPP{}~points relative
to a category with no exposure at all. \StructuralAboveCeiling{} of the
\StructuralHorizons{} horizons imply more than that, including both of the
significant ones. Those estimates are not large, they are impossible: no
allocation of the wage bill produces them. The profile alternates sign between
adjacent horizons as well, which no cumulative response function does, and the
interval at the peak, \StructuralPeakIntervalLower{} to
\StructuralPeakIntervalUpper{}, excludes zero and every attainable value with it.

That diagnosis is one this paper has already applied once. An earlier version of
the regional design returned a peak implying prices rising by more than the
shock that caused them and was rejected on exactly this ground. The same
standard applies here, and the fuller design fails it in the same way.

The falsification test does not rescue the estimates, and it is worth saying why
it cannot. The \StructuralPreTrendLeads{} leads jointly
\StructuralPreTrendRejects{} the null, at $p=\ $\StructuralPreTrendP{}. Passing
a pre-trend test establishes that the leads carry no signal; it says nothing
about whether the contemporaneous coefficient is a magnitude the economics
permits. A design can be clean on timing and still be dividing by a number close
to zero, and the identifying spread of
\StructuralIdentifyingSpread{}~percentage points is that number.

The conclusion is therefore the same as the other two designs reach, arrived at
from the opposite direction. The exposure design fails to reject anything and
cannot bound the coefficient; this one rejects at two horizons and bounds it away
from every value pass-through could take. Neither is evidence about prices in
Portugal. What the fuller design does settle is a question the paper could not
otherwise answer: the binding constraint is not the missing structure. We built
the structure, it absorbs the confounds the other designs are exposed to, and it
still cannot identify a price response, because policy is assigned to nine
regions and no amount of consumption detail adds a tenth.

\subsection{What would change the conclusion}

An earlier version of this section named registering Madeira's missing acts as
the most valuable step available, and estimated that it would roughly quadruple
the identifying events. It has since been done: the acts from 2010 onwards are
registered from primary law, the identifying events rose from three to
\IdentifyingEvents{}, and the estimates became far more plausible without
becoming significant. We record the prediction and its outcome because the
second half is the informative part. More identifying variation of the same kind
improved the point estimates and did not change the conclusion, which is what
one should expect when the binding constraint is that there is one treated
region rather than that there are too few dates.

A second prediction has since been tested and it failed in the more useful
direction. An earlier version argued that building the structural
region-by-category exposure would not help, because the terms it adds are
national. The reasoning was wrong, the design was built, and it does what the
argument for it said it would: it carries region-time and category-time effects
together, so it is exposed to neither the island supply shocks that threaten the
exposure design nor the seasonal cycle that defeats the category one. It still
cannot identify a price response, and it fails in a way that is diagnostic
rather than merely inconclusive, implying differential responses larger than
complete pass-through of the whole minimum-wage cost bill could produce.
Removing the confounds did not help because the confounds were not the
constraint.

That leaves two things, in descending order of value.
Minimum-wage coverage published by NUTS~II region and economic activity jointly
would remove the assumption the exposure measure rests on and, more importantly,
could widen a spread that industry mix alone cannot. And a longer regional price
series at consumption-purpose detail would allow pre-trend testing over a window
long enough to be informative.

Neither would address what both failed designs point at. Policy is assigned to
nine regions and one of them supplies the divergence, so inference clusters on
nine however finely the panel is cut. A third design with more structure would
add cells and not clusters, which is what this one demonstrated.

A fourth would help the exposure design without rescuing it. The published bite
dates from \ExposureBitePeriod{} and so post-dates the 2015 restart of statutory
increases, which is the episode with the most policy variation to exploit; an
earlier monitoring report would make exposure predetermined for it. It would not
widen the \ExposureSpreadPP{}~percentage-point spread, and that, not the dating,
is the binding constraint.
